% =====================================================================
\section{Expected Results and Performance Evaluation}
% =====================================================================

\subsection{Deliverables}

Table~\ref{tab:deliverables} lists the artifacts that will be handed over, with their state on the
proposal date.

\begin{table}[tbp]
\centering
\caption{Deliverables.}
\label{tab:deliverables}
\begin{tabularx}{\textwidth}{@{}L P{38mm} l@{}}
\toprule
Deliverable & Content & State \\
\midrule
Source code repository with README & Model, FastAPI backend, tests & \Done \\
 & React frontend, release tag \code{v1.0} & \Plan \\
Deployment configuration & Docker Compose with PostgreSQL, backup and restore scripts, \code{.env.example} & \Done \\
 & \code{web} service, tunnel token, production settings & \Plan \\
API documentation & OpenAPI 3.1 (Swagger UI, ReDoc), one page per endpoint, database schema & \Done \\
CI/CD workflows & GitHub Actions for pull requests and releases & \Plan \\
Model package and evaluation reports & Frozen checkpoint with manifest; statistics, parity, latency, user study & \Plan \\
User manual and architecture document & Operator guide, administrator guide, design decisions (ADRs) & \Plan \\
Thesis chapters and demo & Chapters 3--5, demo script, backup video & \Plan \\
\bottomrule
\end{tabularx}
\end{table}

\subsection{Verification and Validation Plan}
\label{sec:vv}

Table~\ref{tab:vv} gives the test tiers. Each tier has a tool and a pass criterion that can be checked
from a script in the repository.

\begin{table}[tbp]
\centering
\caption{Verification and validation plan.}
\label{tab:vv}
\begin{tabularx}{\textwidth}{@{}>{\bfseries}P{22mm} L L@{}}
\toprule
Testing tier & Methodology and tools & Success criteria / target metrics \\
\midrule
Unit & pytest (backend), Vitest (frontend) & $\ge 80\%$ backend statement coverage; all tests pass \\
Integration & API against a real PostgreSQL container & Every endpoint returns the documented status code and body \\
Model parity & API vs.\ offline evaluation on 999 validation tracks, greedy and constrained & 0 divergent tracks with the same checkpoint hash \\
Model validation & 3 seeds; McNemar \cite{mcnemar1947} and bootstrap CI \cite{efron1993bootstrap}; sealed test set used once & 95\% CI of the gain over the baseline excludes 0; $p<0.05$ \\
Performance & Load script and Locust on the demo host & p95 $\le 1$\,s per track; 1{,}000 tracks in $\le 20$\,min; reads p95 $\le 100$\,ms during a batch \\
Robustness & Corrupt images, wrong types, too many frames, oversized or broken ZIPs & $\ge 8$ error cases, each with its documented 4xx code; no \code{500} \\
Security & Bandit, pip-audit, npm audit, Trivy, OWASP ZAP baseline \cite{owasp2021} & 0 high or critical findings \\
End to end & 30 cases (20 validation, 10 test tracks) on Chrome Android and Safari iOS & 30/30 pass \\
Usability & 5 users, SUS questionnaire, timed review with and without the app & SUS $\ge 68$; review time per flagged track $-30\%$ \\
\bottomrule
\end{tabularx}
\end{table}

\textbf{Results so far.} The Week~2 backend passes 12 of 12 automated tests. Sending the 999
validation tracks through the API gives the same plates as the offline evaluation for all 999 tracks,
with greedy and with constrained decoding, both natively and in Docker; accuracy through the API is
$79.78\%$, equal to the offline seed-42 run. In a load test with a 40-track batch writing, five clients
posting single tracks and five clients polling, review-queue reads took 8\,ms at p50 and 46\,ms at p95,
batch-progress reads 6\,ms and 20\,ms, with no errors. Single-track latency on CPU is about 270\,ms
natively on the development laptop and 192.75\,ms on the lab server; inside Docker Desktop on macOS it
is 810--910\,ms, so the official measurement will be taken on the Linux demo host in Week~6.
Fig.~\ref{fig:multiseed} shows the multi-seed model results.

\fig{multiseed.pdf}{Exact-match accuracy on the 999 validation tracks. Error bars show the standard
deviation over seeds 42, 100 and 2026.}{fig:multiseed}

On the same tracks, the super-resolved frames reach a PSNR of 16.59\,dB and an SSIM of 0.397, against
15.87\,dB and 0.354 for bicubic upscaling. Across the three seeds, 741 tracks are read correctly every
time and 143 are wrong every time; the second group is the input for the error analysis in Week~5.

\subsection{Success Metrics and SLA Compliance}

Table~\ref{tab:kpi} lists the indicators used to judge the final system, with the value measured so
far where one exists.

\begin{table}[tbp]
\centering
\caption{Key performance indicators.}
\label{tab:kpi}
\begin{tabularx}{\textwidth}{@{}L P{27mm} P{22mm}@{}}
\toprule
KPI & Target & Current \\
\midrule
Exact-match accuracy, validation (3 seeds) & $\ge 79\%$ & 79.95\% $\pm$ 0.15 \\
Exact-match accuracy, sealed test set & reported once with 95\% CI & W6 \\
Precision of auto-accepted tracks & $\ge 95\%$ & W5 \\
Coverage (share auto-accepted) & reported at $\tau^{*}$ & 88.8\% at $\tau=0.90$ (provisional) \\
API vs.\ offline parity & 0 divergent tracks & 0 of 999 \\
Single-track latency on CPU & p95 $\le 1$\,s & $\approx$0.27\,s (native) \\
Batch of 1{,}000 tracks & $\le 20$\,min, 0 errors & $\approx$11\,min (estimate) \\
Availability during demo and user tests & $\ge 99\%$ of health checks OK & W9--12 \\
Data protection & RPO $\le 6$\,h, RTO $< 5$\,min & restore tested on Oct.\ 22 \\
Usability & SUS $\ge 68$ & W10 \\
Review time per flagged track & $-30\%$ vs.\ manual & W10 \\
\bottomrule
\end{tabularx}
\end{table}
